\section*{Chapitre \ref{ch:b2:surfaces} --- Surfaces}

\begin{solution}{exo:b2:surfaces:1}
Avec $\sigma(\theta, r) = (r\cos\theta,\ r\sin\theta,\ r)$~:
\[
\sigma_\theta = (-r\sin\theta,\ r\cos\theta,\ 0),
\qquad
\sigma_r = (\cos\theta,\ \sin\theta,\ 1),
\]
indépendants pour $r > 0$. Le \omterm{def:b2:surfaces:tangent}{plan tangent} en $M_0 =
\sigma(\theta_0, r_0)$ passe par $M_0$ avec les directions
$\sigma_\theta, \sigma_r$. Or $M_0 - O = r_0(\cos\theta_0,
\sin\theta_0, 1) = r_0\,\sigma_r(\theta_0, r_0)$ est lui-même une
direction \omterm{def:b2:curves:arc}{tangente}~: le sommet $O$ est sur le \omterm{def:b2:surfaces:tangent}{plan tangent}. (C'est
le comportement général des cônes~: ils sont réglés par des droites
passant par le sommet, et un \omterm{def:b2:surfaces:tangent}{plan tangent} contient la génératrice
passant par le point de tangence.)
\end{solution}

\begin{solution}{exo:b2:surfaces:2}
Appliquer la~\cref{prop:b2:surfaces:gradient} à $f(x,y,z) =
\frac{x^2}{a^2} + \frac{y^2}{b^2} + \frac{z^2}{c^2}$~: $\nabla
f(x_0,y_0,z_0) = 2\bigl(\frac{x_0}{a^2}, \frac{y_0}{b^2},
\frac{z_0}{c^2}\bigr) \neq 0$ sur la surface. Le \omterm{def:b2:surfaces:tangent}{plan tangent} est
\[
\frac{x_0}{a^2}(x - x_0) + \frac{y_0}{b^2}(y - y_0)
+ \frac{z_0}{c^2}(z - z_0) = 0,
\qquad\text{c.-à-d.}\qquad
\frac{x_0\,x}{a^2} + \frac{y_0\,y}{b^2} + \frac{z_0\,z}{c^2} = 1,
\]
en utilisant que $(x_0, y_0, z_0)$ vérifie l'équation de
l'ellipsoïde --- la règle «~dédoubler les carrés~» qui généralise
$\langle M_0, M \rangle = R^2$ de la sphère.
\end{solution}

\begin{solution}{exo:b2:surfaces:3}
$\sigma_u = (-v\sin u,\ v\cos u,\ a)$ et $\sigma_v = (\cos u,\
\sin u,\ 0)$, donc
\[
E = v^2 + a^2, \qquad F = 0, \qquad G = 1,
\qquad
\sqrt{EG - F^2} = \sqrt{v^2 + a^2} > 0
\]
(l'hélicoïde est régulier partout, y compris sur son axe $v =
0$). \omterm{def:b2:surfaces:area}{Aire} du morceau~:
\[
\mathcal{A} = \int_0^{2\pi}\!\!\int_0^1
\sqrt{v^2 + a^2}\;\dd v\,\dd u
= 2\pi\cdot\frac12\Bigl[v\sqrt{v^2 + a^2}
+ a^2\ln\bigl(v + \sqrt{v^2 + a^2}\bigr)\Bigr]_0^1 ,
\]
soit $\mathcal{A} = \pi\Bigl(\sqrt{1 + a^2} +
a^2\ln\frac{1 + \sqrt{1 + a^2}}{a}\Bigr)$.
\end{solution}

\begin{solution}{exo:b2:surfaces:4}
Dérivées~:
\[
\sigma_\theta = \bigl(-(R + r\cos\psi)\sin\theta,\
(R + r\cos\psi)\cos\theta,\ 0\bigr),
\qquad
\sigma_\psi = \bigl(-r\sin\psi\cos\theta,\
-r\sin\psi\sin\theta,\ r\cos\psi\bigr).
\]
Alors
\[
E = (R + r\cos\psi)^2,
\qquad
F = r\sin\psi\cos\psi\,(R + r\cos\psi)
\bigl(\sin\theta\cos\theta - \sin\theta\cos\theta\bigr) = 0,
\qquad
G = r^2 ,
\]
donc $\sqrt{EG - F^2} = r(R + r\cos\psi) \geq r(R - r) > 0$~:
régulier partout. \omterm{def:b2:surfaces:area}{Aire}~:
\[
\mathcal{A}
= \int_0^{2\pi}\!\!\int_0^{2\pi} r(R + r\cos\psi)
\,\dd\theta\,\dd\psi
= 2\pi r \Bigl(2\pi R + r\int_0^{2\pi}\cos\psi\,\dd\psi\Bigr)
= 4\pi^2 R r ,
\]
le terme en $\cos\psi$ s'intégrant à zéro --- théorème de Pappus~: \omterm{def:b2:surfaces:area}{aire}
$=$ (\omterm{def:b2:curves:length}{longueur} du cercle en rotation) $\times$ (distance parcourue par
son centre).
\end{solution}

\begin{solution}{exo:b2:surfaces:5}
Pour $\sigma(x, y) = (x, y, f(x,y))$~: $\sigma_x = (1, 0, f_x)$,
$\sigma_y = (0, 1, f_y)$, donc $E = 1 + f_x^2$, $F = f_xf_y$, $G = 1
+ f_y^2$ et
\[
EG - F^2 = (1 + f_x^2)(1 + f_y^2) - f_x^2f_y^2
= 1 + f_x^2 + f_y^2 ,
\]
donnant la formule d'\omterm{def:b2:surfaces:area}{aire} annoncée. Pour $f = \frac12(x^2 + y^2)$ sur
le disque unité~: $1 + f_x^2 + f_y^2 = 1 + x^2 + y^2$, et en
coordonnées polaires ($x = \rho\cos\alpha$, $y = \rho\sin\alpha$, jacobien
$\rho$, \cref{ch:b2:multint})~:
\[
\mathcal{A}
= \int_0^{2\pi}\!\!\int_0^1 \sqrt{1 + \rho^2}\;\rho
\,\dd\rho\,\dd\alpha
= 2\pi\Bigl[\tfrac13(1 + \rho^2)^{3/2}\Bigr]_0^1
= \frac{2\pi}{3}\bigl(2\sqrt2 - 1\bigr).
\]
\end{solution}

\begin{solution}{exo:b2:surfaces:6}
D'après le calcul de l'\cref{ex:b2:surfaces:spherearea}, $E =
R^2\cos^2\varphi$, $F = 0$, $G = R^2$, donc par la~\cref{prop:b2:surfaces:curvelength}
\[
L = \int_a^b R\sqrt{\cos^2\varphi(t)\,\theta'(t)^2
+ \varphi'(t)^2}\;\dd t .
\]
Soit les extrémités $(\theta_0, \varphi_1)$ et $(\theta_0,
\varphi_2)$, $\varphi_1 < \varphi_2$. Pour toute courbe joignante,
\[
L \geq \int_a^b R\,\abs{\varphi'(t)}\,\dd t
\geq R\,\Bigl|\int_a^b \varphi'(t)\,\dd t\Bigr|
= R\,(\varphi_2 - \varphi_1),
\]
en laissant tomber le terme positif $\cos^2\varphi\,\theta'^2$ et en utilisant
l'inégalité triangulaire pour les intégrales. L'arc de méridien
$\theta \equiv \theta_0$, $\varphi$ croissant de $\varphi_1$ à
$\varphi_2$, a pour \omterm{def:b2:curves:length}{longueur} exactement $R(\varphi_2 - \varphi_1)$~: il est
le plus court. (Les méridiens sont des grands cercles~; c'est le premier,
élémentaire cas du fait que les géodésiques de la sphère sont les grands
cercles.)
\end{solution}

\begin{solution}{exo:b2:surfaces:7}
Fixons une courbe $\gamma$ tracée sur $S$ et posons $g(t) = \norm{\gamma(t)
- \Omega}^2$. Alors $g'(t) = 2\langle \gamma'(t),\ \gamma(t) -
\Omega\rangle$. La droite normale en $M = \gamma(t)$ passe par
$\Omega$ par hypothèse, donc $\gamma(t) - \Omega$ est un vecteur \emph{normal},
orthogonal au \omterm{def:b2:surfaces:tangent}{plan tangent}, en particulier à la vitesse $\gamma'(t)$
(\cref{prop:b2:surfaces:velocity})~: $g' =
0$, et $g$ est constante le long de toute courbe tracée.

Maintenant l'ensemble $S_c = \{M \in S : \norm{M - \Omega}^2 = c\}$ est
fermé dans $S$~; il est aussi \omterm{def:b2:metric:topology}{ouvert} dans $S$~: au voisinage de chacun de
ses points, $S$ est un graphe régulier, donc tout point voisin de $S$ lui
est joint par une courbe tracée (un segment relevé), le long de laquelle $g$
est constante. Comme $S$ est \omterm{def:b2:metric:connected}{connexe} et $S_c$ non vide pour le bon $c$, $S = S_c
\subseteq$ la sphère de centre $\Omega$ et de rayon $\sqrt c$
(\cref{ch:b2:metric}~: argument de connexité).
\end{solution}

\begin{solution}{exo:b2:surfaces:8}
Écrivons $\Phi(u', v') = (u, v)$. Par la règle de dérivation en chaîne,
\[
\tilde\sigma_{u'} = \frac{\partial u}{\partial u'}\sigma_u
+ \frac{\partial v}{\partial u'}\sigma_v,
\qquad
\tilde\sigma_{v'} = \frac{\partial u}{\partial v'}\sigma_u
+ \frac{\partial v}{\partial v'}\sigma_v ,
\]
dérivées partielles de $\sigma$ évaluées en $\Phi(u',v')$. En développant le
produit vectoriel par bilinéarité et en utilisant $\sigma_u \wedge \sigma_u =
\sigma_v \wedge \sigma_v = 0$, $\sigma_v \wedge \sigma_u =
-\sigma_u \wedge \sigma_v$~:
\[
\tilde\sigma_{u'} \wedge \tilde\sigma_{v'}
= \Bigl(\frac{\partial u}{\partial u'}
\frac{\partial v}{\partial v'}
- \frac{\partial v}{\partial u'}
\frac{\partial u}{\partial v'}\Bigr)\,
\sigma_u \wedge \sigma_v
= \det J_\Phi \cdot (\sigma_u \wedge \sigma_v)\circ\Phi .
\]
En prenant les \omterm{def:b2:nvs:norm}{normes} on obtient l'identité. Puis, par la formule de
changement de variables (\cref{ch:b2:multint}) appliquée à l'application
$\Phi$ sur $K' = \Phi^{-1}(K)$~:
\[
\iint_{K'} \norm{\tilde\sigma_{u'} \wedge \tilde\sigma_{v'}}
\,\dd u'\dd v'
= \iint_{K'} \norm{\sigma_u \wedge \sigma_v}\circ\Phi\;
\abs{\det J_\Phi}\,\dd u'\dd v'
= \iint_K \norm{\sigma_u \wedge \sigma_v}\,\dd u\,\dd v :
\]
les deux paramétrages attribuent la même \omterm{def:b2:surfaces:area}{aire} au même morceau de
surface.
\end{solution}

\begin{solution}{exo:b2:surfaces:9}
Dans la carte sphérique, $z = R\sin\varphi$, et la zone
correspond à $\varphi_1 \leq \varphi \leq \varphi_2$ avec
$z_i = R\sin\varphi_i$. Avec l'élément d'\omterm{def:b2:surfaces:area}{aire}
$R^2\cos\varphi\,\dd\theta\,\dd\varphi$
(\cref{ex:b2:surfaces:spherearea})~:
\[
\mathcal A = \int_{\varphi_1}^{\varphi_2}\!\!\int_0^{2\pi}
R^2\cos\varphi\,\dd\theta\,\dd\varphi
= 2\pi R^2(\sin\varphi_2 - \sin\varphi_1)
= 2\pi R\,(z_2 - z_1) .
\]
Le résultat ne dépend que de la hauteur $z_2 - z_1$~: des tranches d'épaisseur
égale portent des \omterm{def:b2:surfaces:area}{aires} égales, qu'elles soient coupées à l'équateur
ou au pôle --- le théorème de la boîte à chapeau d'Archimède, et la raison
pour laquelle l'\omterm{def:b2:surfaces:area}{aire} latérale du cylindre circonscrit ($2\pi R \cdot
2R = 4\pi R^2$) égale l'\omterm{def:b2:surfaces:area}{aire} de la sphère.
\end{solution}

\begin{solution}{exo:b2:surfaces:10}
$\sigma_\theta = (-r\sin\theta,\ r\cos\theta,\ 0)$ et
$\sigma_z = (r'\cos\theta,\ r'\sin\theta,\ 1)$, donc
\[
\sigma_\theta \wedge \sigma_z
= (r\cos\theta,\ r\sin\theta,\ -r\,r') ,
\]
un vecteur normal en $M = (r\cos\theta, r\sin\theta, z)$. La
droite normale est
\[
t \mapsto \bigl(r(1 + t)\cos\theta,\ r(1 + t)\sin\theta,\
z - t\,r\,r'\bigr),
\]
qui en $t = -1$ atteint $(0,\ 0,\ z + r(z)\,r'(z))$~: toute
droite normale rencontre l'axe, à la hauteur $z + rr'$. (C'est
la raison tridimensionnelle pour laquelle la symétrie de rotation survit dans
le champ normal.)
\end{solution}

\begin{solution}{exo:b2:surfaces:11}
$\sigma_u = (-\sin u, \cos u, 0)$, $\sigma_v = (0, 0, 1)$~: $E
= 1$, $F = 0$, $G = 1$. Par la~\cref{prop:b2:surfaces:curvelength} la \omterm{def:b2:curves:length}{longueur} d'une courbe
tracée est $\int\sqrt{u'^2 + v'^2}\,\dd t$ --- la \omterm{def:b2:curves:length}{longueur} de son
ombre paramétrique $(u(t), v(t))$ dans le plan~: la carte est une
isométrie locale (le déroulement du cylindre). L'hélice
$t \mapsto (\cos t, \sin t, ct)$, $t \in \intcc0{2\pi}$, a pour
ombre le segment de $(0,0)$ à $(2\pi, 2\pi c)$, de \omterm{def:b2:curves:length}{longueur}
$2\pi\sqrt{1 + c^2}$. Toute courbe tracée de $(1,0,0)$ à $(1,
0, 2\pi c)$ faisant un tour \omterm{def:b2:metric:complete}{complet} a une ombre \omterm{def:b2:metric:continuity}{continue}
joignant $(0, 0)$ à $(2\pi, 2\pi c)$, de \omterm{def:b2:curves:length}{longueur} plane $\geq$
le segment droit~; puisque les \omterm{def:b2:curves:length}{longueurs} coïncident, l'hélice est
la plus courte.
\end{solution}

\begin{solution}{exo:b2:surfaces:12}
La fonction $g(M) = \norm{M}^2$ est \omterm{def:b2:metric:continuity}{continue} sur le \omterm{def:b2:metric:compact}{compact}
$S$, donc elle atteint son maximum en un certain $M_0$. Pour toute courbe
$\gamma$ tracée sur $S$ avec $\gamma(0) = M_0$, la fonction $t
\mapsto \norm{\gamma(t)}^2$ a un maximum en $t = 0$, donc sa
dérivée $2\langle\gamma'(0), M_0\rangle$ s'annule~: $M_0$ est
orthogonal à tout vecteur tangent, c.-à-d.\ normal à $S$ en
$M_0$. Comme $\nabla f(M_0) \neq 0$ dirige aussi la droite normale
(\cref{prop:b2:surfaces:gradient}), $\nabla f(M_0)$ et
$\vect{OM_0}$ sont colinéaires. Pour l'ellipsoïde, la radialité
signifie
\[
\Bigl(\frac{x_0}{a^2}, \frac{y_0}{b^2},
\frac{z_0}{c^2}\Bigr) = \mu\,(x_0, y_0, z_0) :
\]
chaque coordonnée vérifie $x_0(\frac1{a^2} - \mu) = 0$, etc.~;
puisque $a^{-2}, b^{-2}, c^{-2}$ sont distincts, au plus une
coordonnée est non nulle, et les solutions sur la surface sont
les six extrémités des axes $(\pm a, 0, 0)$, $(0, \pm b, 0)$, $(0,
0, \pm c)$. Les points les plus éloignés sont $(\pm a, 0, 0)$, à
distance $a = \max(a,b,c)$.
\end{solution}

\begin{solution}{pb:b2:surfaces:1}
\textbf{1.} En développant, et en utilisant la symétrie de $A$
($t^{\mathsf T}\!APY = (At)^{\mathsf T}PY$)~:
\[
q(PY + t) = Y^{\mathsf T}P^{\mathsf T}\!APY +
2\,(At + b)^{\mathsf T}PY + \bigl(t^{\mathsf T}\!At +
2b^{\mathsf T}t + c\bigr),
\]
qui est le triplet affiché. $A' = P^{\mathsf T}\!AP =
P^{-1}AP$ est semblable à $A$~: même \omterm{def:b2:reduction:charpoly}{polynôme caractéristique},
même \omterm{def:b2:reduction:eigen}{spectre}, même rang, même signature. Enfin $\{q = 0\} = \{sq = 0\}$
pour $s \neq 0$, donc seule l'équation \emph{à un scalaire près}
est attachée à l'ensemble~; multiplier par $s$ multiplie toutes les \omterm{def:b2:reduction:eigen}{valeurs
propres} par $s$.

\textbf{2.} Le théorème spectral fournit $P \in O(3)$ avec
$P^{\mathsf T}\!AP = \operatorname{diag}(\lambda_1, \lambda_2,
\lambda_3)$~; la question 1 avec $t = 0$ transforme l'équation en
$\sum\lambda_iy_i^2 + 2\sum\beta_iy_i + c = 0$, où $\beta =
P^{\mathsf T}b$.

\textbf{3.} Pour $\lambda_i \neq 0$~: $\lambda_iy_i^2 +
2\beta_iy_i = \lambda_i\bigl(y_i +
\frac{\beta_i}{\lambda_i}\bigr)^2 -
\frac{\beta_i^2}{\lambda_i}$~; la translation $z_i = y_i +
\beta_i/\lambda_i$ (et $z_i = y_i$ pour $i > r$) donne
\[
\sum_{i \leq r}\lambda_iz_i^2 + 2\sum_{i > r}\beta_iz_i +
c'' = 0, \qquad
c'' = c - \sum_{i\leq r}\frac{\beta_i^2}{\lambda_i}.
\]

\textbf{4.} $q(2\Omega - X) = (2\Omega - X)^{\mathsf T}
A(2\Omega - X) + 2b^{\mathsf T}(2\Omega - X) + c$~; en développant
et en soustrayant $q(X)$, les termes quadratiques s'annulent et
\[
q(2\Omega - X) - q(X) = -4\,\langle A\Omega + b,\ X\rangle +
4\,\langle A\Omega + b,\ \Omega\rangle .
\]
Si $A\Omega + b = 0$ ceci s'annule identiquement~: la symétrie
centrale préserve $q$, donc la \omterm{pb:b2:surfaces:1}{quadrique}. Réciproquement,
«~$q(2\Omega - X) = q(X)$ pour tout $X$~» dit que la fonction \omterm{def:b2:affine:subspace}{affine}
ci-dessus s'annule sur tout $\R^3$, ce qui force sa partie linéaire
$A\Omega + b$ à être nulle. Donc centres $=$ solutions de
$A\Omega = -b$~: un ensemble non vide ssi $b \in \operatorname{im}A$
(un sous-espace \omterm{def:b2:affine:subspace}{affine} dirigé par $\ker A$), et un point unique
ssi $A$ est inversible.

\textbf{5.} Signature $(3,0)$ (tous les $\lambda_i > 0$)~: $\delta >
0$ donne $\frac{x^2}{a^2} + \frac{y^2}{b^2} + \frac{z^2}{c^2}
= 1$ avec $a = \sqrt{\delta/\lambda_1}$, etc.\ --- un
ellipsoïde~; $\delta = 0$~: le point unique $O$~; $\delta < 0$~:
vide. Signature $(2,1)$ ($\lambda_1, \lambda_2 > 0 >
\lambda_3$)~: $\delta > 0$~: $\frac{x^2}{a^2} + \frac{y^2}{b^2}
- \frac{z^2}{c^2} = 1$, l'hyperboloïde à une nappe~; $\delta =
0$~: le cône $\frac{x^2}{a^2} + \frac{y^2}{b^2} =
\frac{z^2}{c^2}$~; $\delta < 0$~: $\frac{z^2}{c^2} -
\frac{x^2}{a^2} - \frac{y^2}{b^2} = 1$, l'hyperboloïde à deux
nappes ($\abs z \geq c$~: deux composantes).

\textbf{6.} La partie quadratique a pour matrice $A$ avec diagonale
$1$ et hors-diagonale $2$~: $A = 2J - I$. Puisque $J$ a pour \omterm{def:b2:reduction:eigen}{spectre}
$\{3, 0, 0\}$ (vecteur propre $(1,1,1)$ pour $3$), $A$ a pour
\omterm{def:b2:reduction:eigen}{spectre} $\{5, -1, -1\}$, la \omterm{def:b2:reduction:eigen}{valeur propre} $5$ portée par
$\R(1,1,1)$. Dans les coordonnées tournées~: $5u^2 - v^2 - w^2 =
1$, c.-à-d.\ $\frac{u^2}{1/5} - v^2 - w^2 = 1$~: un hyperboloïde à
deux nappes, de révolution (\omterm{def:b2:reduction:eigen}{valeurs propres} égales $-1$) autour de
l'axe $\R(1,1,1)$.

\textbf{7.} La surface est $\bigl(\frac xa - \frac
zc\bigr)\bigl(\frac xa + \frac zc\bigr) = \bigl(1 - \frac
yb\bigr)\bigl(1 + \frac yb\bigr)$. Pour $(\lambda : \mu) \neq
(0:0)$ définir la droite
\[
D_{\lambda:\mu} :\quad
\lambda\Bigl(\frac xa - \frac zc\Bigr) = \mu\Bigl(1 - \frac
yb\Bigr), \qquad
\mu\Bigl(\frac xa + \frac zc\Bigr) = \lambda\Bigl(1 + \frac
yb\Bigr)
\]
(deux équations \omterm{def:b2:affine:subspace}{affines} indépendantes~: une droite). En multipliant les
deux équations on voit que tout point de $D_{\lambda:\mu}$ est
sur la surface lorsque $\lambda\mu \neq 0$~; les cas $\lambda = 0$
ou $\mu = 0$ se vérifient directement (par exemple $\lambda = 0$~: $y =
b$, $\frac xa = -\frac zc$, ce qui satisfait l'équation). La
seconde famille $D'_{\lambda:\mu}$ échange les deux facteurs de
droite.

\textbf{8.} Fixons $M$ sur la surface. Les conditions pour $M \in
D_{\lambda:\mu}$ forment un système linéaire homogène $2\times2$
en $(\lambda, \mu)$ dont le \omterm{def:b2:linalg:det}{déterminant} est
\[
\Bigl(\frac{x^2}{a^2} - \frac{z^2}{c^2}\Bigr) - \Bigl(1 -
\frac{y^2}{b^2}\Bigr) = 0
\]
précisément parce que $M$ est sur la \omterm{pb:b2:surfaces:1}{quadrique}~: une solution non triviale
$(\lambda : \mu)$ existe. La matrice des coefficients n'est
jamais nulle (cela forcerait $1 - \frac yb = 1 + \frac yb =
0$), donc son rang est $1$ et la solution est unique à un facteur près~:
exactement une droite de la famille passe par $M$.
Il en va de même pour la seconde famille, et les deux droites sont
distinctes (en $(a, 0, 0)$ ce sont $\{x = a,\ \frac zc = \frac
yb\}$ et $\{x = a,\ \frac zc = -\frac yb\}$)~: l'hyperboloïde à une
nappe est doublement réglé.

\textbf{9.} Soit $M = (x, y, z)$ sur l'hyperboloïde, $\rho^2 =
\frac{x^2}{a^2} + \frac{y^2}{b^2} = 1 + \frac{z^2}{c^2} \geq
1$. Le point $N = (x,\ y,\ \varepsilon c\rho)$ avec
$\varepsilon$ le signe de $z$ vérifie $\frac{x^2}{a^2} +
\frac{y^2}{b^2} - \frac{(c\rho)^2}{c^2} = 0$~: $N \in C$, et
\[
d(M, C) \leq \abs{z - \varepsilon c\rho}
= c\bigl(\rho - \sqrt{\rho^2 - 1}\bigr)
= \frac{c}{\rho + \sqrt{\rho^2 - 1}} .
\]
Si $\norm M \to \infty$ alors $\rho \to \infty$ (les trois
coordonnées sont majorées par des multiples de $\rho$), donc $d(M, C)
\to 0$. La section $x = a$~: $\frac{y^2}{b^2} -
\frac{z^2}{c^2} = 0$, la paire de droites sécantes de la question
8 --- et de même en $x = -a$.

\textbf{10.} Avec $\lambda_3 = 0$, la question 3 laisse
$\lambda_1z_1^2 + \lambda_2z_2^2 + 2\beta_3z_3 + c'' = 0$. Dans
la base propre, $\operatorname{im}A = \operatorname{Vect}(e_1,
e_2)$, donc par la question 4 les centres existent ssi $\beta_3 = 0$. Si
$\beta_3 \neq 0$~: la translation $z_3 \mapsto z_3 -
c''/(2\beta_3)$ supprime la constante, laissant $\lambda_1z_1^2
+ \lambda_2z_2^2 + 2\beta_3z_3 = 0$, c.-à-d.\ $z_3 = px^2 +
qy^2$ après renommage~: un \emph{paraboloïde elliptique} si
$\lambda_1\lambda_2 > 0$, un \emph{paraboloïde hyperbolique} si
$\lambda_1\lambda_2 < 0$ --- et effectivement pas de centre. Si
$\beta_3 = 0$~: l'équation $\lambda_1z_1^2 + \lambda_2z_2^2 +
c'' = 0$ ne fait pas intervenir $z_3$~: la \omterm{pb:b2:surfaces:1}{quadrique} est un cylindre
sur la conique plane correspondante --- cylindre elliptique,
droite, ou ensemble vide lorsque $\lambda_1\lambda_2 > 0$~; cylindre
hyperbolique ou paire de plans sécants lorsque
$\lambda_1\lambda_2 < 0$ --- avec toute une droite de centres
$\{(z_1^*, z_2^*)\} \times \R$.

\textbf{11.} $xy - z = 0$. En substituant $x = \frac{u +
v}{\sqrt2}$, $y = \frac{u - v}{\sqrt2}$ (rotation de $\pi/4$)~:
$xy = \frac{u^2 - v^2}2$, donc l'équation devient $z =
\frac12(u^2 - v^2)$~: un paraboloïde hyperbolique --- la selle
de la figure, au facteur $\frac12$ près.

\textbf{12.} La droite $\{x = x_0,\ z = x_0y\}$ (paramétrée
par $y$) est clairement sur $z = xy$, tout comme $\{y = y_0,\ z =
xy_0\}$~; par $(x_0, y_0, x_0y_0)$ passent les deux.
Unicité~: si $t \mapsto (x_0 + tv_1, y_0 + tv_2, z_0 +
tv_3)$ reste sur la surface, le coefficient de $t^2$ dans $(x_0
+ tv_1)(y_0 + tv_2) - z_0 - tv_3$ donne $v_1v_2 = 0$, donc
$v_1 = 0$ ou $v_2 = 0$, tombant dans l'une des deux familles~:
une droite de chaque par chaque point --- la seconde \omterm{pb:b2:surfaces:1}{quadrique} doublement
réglée.

\textbf{13.} En complétant les carrés~: $(x - 1)^2 + (y + 2)^2 = 2$,
sans condition sur $z$~: un cylindre circulaire droit de rayon
$\sqrt2$ et d'axe la droite verticale $\{(1, -2, z) : z \in
\R\}$ --- une droite de centres, comme le prédit la question 10.

\textbf{14.} Avec $\lambda_2 = \lambda_3 = 0$ l'équation réduite
est $\lambda_1z_1^2 + 2\beta_2z_2 + 2\beta_3z_3 + c''
= 0$. Une rotation du plan noyau $(z_2, z_3)$ aligne la
forme linéaire~: $2\beta_2z_2 + 2\beta_3z_3 = 2\beta w$ avec
$\beta = \sqrt{\beta_2^2 + \beta_3^2}$. Si $\beta \neq 0$,
translater $w$ pour absorber $c''$~: $\lambda_1z_1^2 + 2\beta w =
0$, un \emph{cylindre parabolique}~; si $\beta = 0$~:
$\lambda_1z_1^2 = -c''$ donne deux plans parallèles ($c''
\lambda_1 < 0$), un plan double ($c'' = 0$), ou l'ensemble
vide. Pour $(x + y)^2 = z$~: avec $u = \frac{x + y}{\sqrt2}$ l'équation
s'écrit $z = 2u^2$~: un cylindre parabolique, invariant
par les translations le long de $(1, -1, 0)$.

\textbf{15.} Toute \omterm{pb:b2:surfaces:1}{quadrique} est envoyée par une rotation plus des
translations sur l'une de~: (rang 3) ellipsoïde, point, ensemble
vide, hyperboloïde à une nappe, cône, hyperboloïde à deux nappes~;
(rang 2) paraboloïde elliptique, paraboloïde hyperbolique, cylindre
elliptique, droite, ensemble vide, cylindre hyperbolique, paire de plans
sécants~; (rang 1) cylindre parabolique, paire de plans parallèles, plan
double, ensemble vide. En identifiant les trois variantes vides comme des
\emph{types \omterm{def:b2:affine:subspace}{affines}} distincts d'équations, le compte est dix-sept~; parmi
elles neuf sont d'honnêtes surfaces~: ellipsoïde, les deux hyperboloïdes,
le cône, les deux paraboloïdes et les trois cylindres.

\textbf{16.} Algorithme. (i) Lire $(A, b, c)$~; calculer le
\omterm{def:b2:reduction:charpoly}{polynôme caractéristique} de $A$, ses \omterm{def:b2:reduction:eigen}{valeurs propres} (réelles, par
le théorème spectral) et $r = \operatorname{rank}A$. (ii)
Résoudre $A\Omega = -b$ (élimination de Gauss)~: résoluble ou non
--- centres ou non. (iii) Si résoluble, translater vers un centre~:
l'équation devient $\sum\lambda_iz_i^2 + c'' = 0$ avec $c''
= c + \langle b, \Omega\rangle$~; classer selon $r$, la signature
et le signe de $c''$ en utilisant les questions 5, 10, 14. (iv) Si non
résoluble ($r \leq 2$), tourner et réduire comme aux questions 10
et 14~: paraboloïde ($r = 2$) ou cylindre parabolique ($r = 1$),
elliptique/hyperbolique selon le signe de
$\lambda_1\lambda_2$. Chaque étape est un calcul fini~: racines
d'un cubique à racines réelles, rangs, systèmes linéaires.

\textbf{17.} $A$ a diagonale $1$, hors-diagonale $-1$~: $A = 2I
- J$, \omterm{def:b2:reduction:eigen}{spectre} $\{2 - 3,\ 2,\ 2\} = \{-1, 2, 2\}$ avec $-1$
sur $\R(1,1,1)$. Signature $(2,1)$, $b = 0$, membre de droite
$\delta = 1 > 0$~: $2u^2 + 2v^2 - w^2 = 1$, un hyperboloïde à une
nappe de révolution autour de l'axe $\R(1,1,1)$.

\textbf{18.} Compléter les carrés~: $(x-1)^2 + (y+2)^2 - (z-1)^2 +
(-1 - 4 + 1 + 4) = 0$, c.-à-d.
\[
(x-1)^2 + (y+2)^2 = (z-1)^2 :
\]
la constante s'est annulée --- un \emph{cône} circulaire droit avec
sommet (et unique centre) $(1, -2, 1)$ et axe parallèle à
$Oz$.

\textbf{19.} La partie quadratique $x^2 + 4xy + y^2$ a pour matrice
$\begin{psmallmatrix}1 & 2\\ 2 & 1\end{psmallmatrix}$ (dans le
plan $xy$), \omterm{def:b2:reduction:eigen}{valeurs propres} $3$ (sur $(1,1)$) et $-1$ (sur
$(1,-1)$)~: avec $u = \frac{x+y}{\sqrt2}$, $v =
\frac{x-y}{\sqrt2}$ elle vaut $3u^2 - v^2$, et la \omterm{pb:b2:surfaces:1}{quadrique}
est
\[
z = \tfrac32u^2 - \tfrac12v^2 :
\]
un paraboloïde hyperbolique ($A$ a rang $2$ et $b$ a une
composante le long de $\ker A = \R e_z$~: pas de centre).

\textbf{20.} Un déplacement transforme les données de l'équation par
$A \mapsto P^{\mathsf T}\!AP$ (mêmes \omterm{def:b2:reduction:eigen}{valeurs propres}) et les
formes normales n'ont aucune liberté résiduelle sauf permuter
les coordonnées et multiplier toute l'équation par un scalaire
($> 0$ pour préserver l'écriture)~: deux formes normales à centre
coïncident à isométrie près ssi les listes de coefficients concordent à
permutation et à un facteur positif commun près --- pour l'ellipsoïde,
ssi les demi-axes $(a, b, c)$ concordent. Affinement, on peut aussi
mettre à l'échelle chaque coordonnée séparément ($z_i \mapsto z_i\sqrt{\abs
{\lambda_i}}$), ce qui efface les \omterm{def:b2:reduction:eigen}{valeurs propres} et ne laisse que
leurs signes~: tout ellipsoïde devient $u^2 + v^2 + w^2 = 1$,
la sphère. Le théorème d'inertie de Sylvester
(\cref{thm:b2:quadratic:sylvester}) garantit que la signature
survit à tout changement linéaire inversible~: les types \omterm{def:b2:affine:subspace}{affines} de
la question 15 sont véritablement distincts.

\textbf{21.} Paramétrer le plan affinement~: $M = P + su +
tv$. Alors $q(P + su + tv)$ est un polynôme de degré $\leq 2$
en $(s, t)$ (développer la \omterm{def:b2:quadratic:def}{forme quadratique} par bilinéarité), donc la
section $\{q = 0\}$ est une conique du plan, éventuellement
dégénérée. Pour $z = xy$ et le plan $z = c$~: $xy = c$, une
hyperbole pour $c \neq 0$, et pour $c = 0$ les deux droites
coordonnées --- la paire de génératrices passant par l'origine.

\textbf{22.} \emph{Ellipsoïde}~: borné, ne contient aucune droite.
\emph{Hyperboloïde à deux nappes} $\frac{z^2}{c^2} -
\frac{x^2}{a^2} - \frac{y^2}{b^2} = 1$~: restreindre à $p + tv$~;
le coefficient de $t^2$ $\frac{v_3^2}{c^2} - \frac{v_1^2}{a^2}
- \frac{v_2^2}{b^2}$ doit s'annuler, donc $v_3 \neq 0$ (sinon $v =
0$)~; le coefficient de $t$ donne $\frac{p_3v_3}{c^2} =
\frac{p_1v_1}{a^2} + \frac{p_2v_2}{b^2}$, et Cauchy--Schwarz
donne
\[
\frac{p_3^2}{c^2}
= \frac{c^2}{v_3^2}\Bigl(\frac{p_1v_1}{a^2} +
\frac{p_2v_2}{b^2}\Bigr)^2
\leq \frac{c^2}{v_3^2}\Bigl(\frac{p_1^2}{a^2} +
\frac{p_2^2}{b^2}\Bigr)\frac{v_3^2}{c^2}
= \frac{p_1^2}{a^2} + \frac{p_2^2}{b^2},
\]
donc le terme constant est $\leq 0 \neq 1$~: aucune droite.
\emph{Paraboloïde elliptique} $z = \frac{x^2}{a^2} +
\frac{y^2}{b^2}$~: le coefficient de $t^2$ force $v_1 = v_2 =
0$, alors l'équation est linéaire non constante en $t$~: aucune droite.
\emph{Cône} $x^2 + y^2 = z^2$~: une droite dessus vérifie $q(p) =
q(v) = B(p, v) = 0$ pour la forme de Lorentz~; l'égalité dans le
Cauchy--Schwarz plan $\abs{p_1v_1 + p_2v_2} = \abs{p_3v_3} =
\sqrt{(p_1^2 + p_2^2)(v_1^2 + v_2^2)}$ force $(p_1, p_2)
\parallel (v_1, v_2)$ et alors $p = kv$~: toutes les droites passent
par le sommet --- une famille. \emph{Cylindres}~: pour les
cylindres elliptique et parabolique le coefficient de $t^2$ force
$v_1 = v_2 = 0$ (seulement les génératrices)~; pour le cylindre hyperbolique
$\frac{x^2}{a^2} - \frac{y^2}{b^2} = 1$, $\frac{v_1}a = \pm
\frac{v_2}b$ avec $v_2 \neq 0$ mène via le coefficient de $t$
à $\frac{p_1^2}{a^2} = \frac{p_2^2}{b^2}$, contredisant le
terme constant $1$~: à nouveau seulement les génératrices verticales. Donc les
surfaces \omterm{pb:b2:surfaces:1}{quadriques} doublement réglées sont exactement l'hyperboloïde à une
nappe et le paraboloïde hyperbolique.

\textbf{23.} Gauss~: $x^2 - 2xy - 2zx = (x - y - z)^2 - y^2 -
z^2 - 2yz$, donc
\[
q = (x - y - z)^2 - 4yz = (x - y - z)^2 + (y - z)^2 - (y +
z)^2 :
\]
trois carrés indépendants de signes $(+, +, -)$ --- signature
$(2, 1)$, en accord avec la question 17, sans aucun calcul de \omterm{def:b2:reduction:eigen}{valeur
propre}. Gauss perd les données métriques~: les nouvelles coordonnées
ne sont pas orthonormées, donc les \omterm{def:b2:reduction:eigen}{valeurs propres} (la forme de
l'hyperboloïde, ses axes et leurs \omterm{def:b2:curves:length}{longueurs}) sont perdues~; seul le
type \omterm{def:b2:affine:subspace}{affine} subsiste.

\textbf{24.} Soient $p$, $n$, $z$ les nombres de \omterm{def:b2:reduction:eigen}{valeurs propres} positives,
négatives et nulles, $p + n + z = 3$. La règle de Descartes
majore $p$ par le nombre $V$ de changements de signe de
$\chi_A$, et $n$ par le nombre $V'$ de changements de signe de
$\chi_A(-\lambda)$~; de plus chaque paire de coefficients consécutifs
non nuls produit un changement dans exactement l'un des deux
polynômes, donc $V + V' \leq 3 - z$ (les racines nulles sont visibles
comme des coefficients de queue nuls). Alors $p + n = 3 - z \geq
V + V' \geq p + n$~: égalité, donc $p = V$ exactement --- les
signes des \omterm{def:b2:reduction:eigen}{valeurs propres} se lisent. Pour $\chi_A(\lambda)
= \lambda^3 - 3\lambda^2 - 9\lambda - 5$~: signes $+,-,-,-$
donnent $V = 1$, et $\chi_A(-\lambda) = -\lambda^3 - 3\lambda^2
+ 9\lambda - 5$ a pour signes $-,-,+,-$~: $V' = 2$. Signature $(1,
2)$ --- cohérent avec la factorisation exacte $\chi_A =
(\lambda - 5)(\lambda + 1)^2$ de la question 6.

\textbf{25.} Algorithme~: diagonaliser la partie quadratique
de manière orthonormée (théorème spectral~: c'est la seule
étape analytiquement profonde, et c'est ce qui rend la
classification \emph{euclidienne})~; résoudre $A\Omega = -b$ pour
décider entre types à centre et paraboliques et pour translater
la partie linéaire là où c'est possible (géométrie \omterm{def:b2:affine:subspace}{affine})~; lire le
type à partir du rang, de la signature et de la constante (Sylvester
garantit que ce sont des invariants)~; quand seul le type \omterm{def:b2:affine:subspace}{affine}
importe, la \omterm{thm:b2:quadratic:gauss}{réduction de Gauss} remplace le théorème spectral au
prix de l'information métrique. Dans $\R^2$ la même
machine classe les coniques~: ellipse, hyperbole, parabole, plus
des paires de droites, une droite, un point et des ensembles vides. Dans $\R^n$
rien ne change sinon la comptabilité~: les types sont indexés par la
signature de $A$, la position de $b$ relativement à
$\operatorname{im}A$, et une constante --- avec la
matrice bordée de dimension $(n{+}2)$ de $(A, b, c)$
fournissant un invariant \omterm{def:b2:metric:compact}{compact}.
\end{solution}
